\documentclass[10pt]{article}
\usepackage[T1]{fontenc}
\usepackage[utf8]{inputenc}
\usepackage[french]{babel}
\usepackage{newpxtext,newpxmath}
\usepackage[a4paper,top=1.5cm,bottom=1.35cm,left=1.6cm,right=1.6cm]{geometry}
\usepackage{microtype}
\usepackage{xcolor}
\usepackage{titlesec}
\usepackage[most]{tcolorbox}
\usepackage{enumitem}
\usepackage{array}
\usepackage{booktabs}
\usepackage{fancyhdr}

\definecolor{navy}{HTML}{1F3A5F}
\definecolor{grey}{HTML}{54606E}
\definecolor{green}{HTML}{2C6E49}
\definecolor{amber}{HTML}{9C6A16}
\definecolor{red}{HTML}{A6322E}
\definecolor{plum}{HTML}{7A4EA8}
\definecolor{boxgray}{HTML}{F3F5F8}
\definecolor{boxred}{HTML}{FCF2F1}
\definecolor{boxgreen}{HTML}{F1F7F3}
\definecolor{boxplum}{HTML}{F7F3FB}
\definecolor{boxcode}{HTML}{F7F8FA}
\definecolor{boxedge}{HTML}{D9E0EA}

\setlength{\parindent}{0pt}
\setlength{\parskip}{2.6pt}
\linespread{1.0}

\pagestyle{fancy}\fancyhf{}
\lhead{\small\color{grey}Scella · architecture}
\rhead{\small\color{grey}M1 --- Le garde unique --- 20 août 2026}
\cfoot{\small\color{grey}\thepage}
\renewcommand{\headrulewidth}{0.3pt}

\titleformat{\section}{\color{navy}\large\bfseries}{}{0pt}{}
\titlespacing{\section}{0pt}{11pt}{3pt}
\titleformat{\subsection}{\color{navy}\bfseries}{}{0pt}{}
\titlespacing{\subsection}{0pt}{7pt}{1.5pt}

\newtcolorbox{these}{colback=boxgray,colframe=boxedge,boxrule=0.5pt,arc=2.5pt,
  left=8pt,right=8pt,top=4.5pt,bottom=4.5pt}
\newtcolorbox{alerte}{colback=boxred,colframe=red!55,boxrule=0.8pt,arc=2.5pt,
  left=8pt,right=8pt,top=4.5pt,bottom=4.5pt}
\newtcolorbox{acquis}{colback=boxgreen,colframe=green!45,boxrule=0.8pt,arc=2.5pt,
  left=8pt,right=8pt,top=4.5pt,bottom=4.5pt}
\newtcolorbox{piste}{colback=boxplum,colframe=plum!45,boxrule=0.8pt,arc=2.5pt,
  left=8pt,right=8pt,top=4.5pt,bottom=4.5pt}
\newtcolorbox{extrait}{colback=boxcode,colframe=boxedge,boxrule=0.5pt,arc=1.5pt,
  left=7pt,right=7pt,top=3.5pt,bottom=3.5pt}

\newcommand{\best}[1]{\textcolor{green}{\textbf{#1}}}
\newcommand{\pire}[1]{\textcolor{red}{\textbf{#1}}}
\newcommand{\sz}[1]{\textcolor{grey}{\small\textsf{#1}}}
\newcommand{\code}[1]{\texttt{\small #1}}
\newcommand{\fl}{$\rightarrow$}

\begin{document}

\begin{center}
{\Huge\bfseries\color{navy} Scella}\\[4pt]
{\large\color{grey}\itshape M1 --- Le garde unique · architecture}\\[4pt]
{\small\color{grey} 20 août 2026 \; · \; second module du socle \; · \; indépendant de M0}
\end{center}
\vspace{2pt}

\begin{these}
\textbf{Le module ne s'appelle pas comme ce qu'il fait.} La roadmap demande d'unifier \emph{un
garde} --- la fonction \code{is\_local\_endpoint}, dupliquée en deux exemplaires. L'inventaire dit
autre chose : l'unité qui compte n'est pas le garde mais l'\textbf{analyseur d'autorité} qui le
nourrit, il en existe \pire{cinq} et non deux, et trois d'entre eux prennent une décision de
sécurité. Le critère de fin écrit dans la roadmap peut être atteint \emph{sans fermer la faille pour
laquelle il a été écrit}. M1 doit donc commencer par corriger sa propre définition.
\end{these}

\section{1 · Ce que la roadmap demande, et ce que le dépôt contient}

La roadmap modulaire fixe M1 en une ligne --- deux jours, aucune dépendance, faisable n'importe
quand --- et donne la preuve du trou sous la même forme, exécutable :

\begin{extrait}
\code{\sz{fini quand :}\ \ grep -rln 'fn is\_local\_endpoint' \ \fl{} \textbf{un} fichier} \\
\code{\sz{aujourd'hui :}\ \ \ \fl{} provenance-capture/src/guard.rs} \\
\code{\sz{\ \ \ \ \ \ \ \ \ \ \ \ \ \ }\ \ \fl{} efe-swarm-sdk/src/backend/mod.rs}
\end{extrait}

C'est exact, et c'est insuffisant. La commande compte \emph{un nom de fonction}. La classe de
vulnérabilité, elle, s'appelle \textbf{différentiel d'analyseur} : elle survient dès que deux morceaux
de code déduisent un hôte d'une même chaîne et n'obtiennent pas le même résultat, ou n'obtiennent pas
celui que le client HTTP contactera réellement. Ce qu'il faut compter, ce sont donc les
\emph{analyseurs}, pas les gardes. Le dépôt en contient cinq.

\begin{center}\footnotesize
\renewcommand{\arraystretch}{1.25}
\begin{tabular}{@{}>{\raggedright\arraybackslash}p{4.4cm} >{\raggedright\arraybackslash}p{4.3cm} >{\raggedright\arraybackslash}p{6.6cm}@{}}
\toprule
\textbf{Emplacement} & \textbf{Ce qu'il décide} & \textbf{Durcissement} \\
\midrule
\code{provenance-capture/} \code{src/guard.rs:11} &
\textcolor{navy}{\textbf{référence}} --- refus d'egress, \code{assert\_local} &
\best{durci} : coupe \code{/ ? \#}, \emph{fail-close} sur \code{\textbackslash}\,/espace/contrôle, \code{userinfo} au \emph{premier} \code{@}, 4 octets numériques exigés \\
\code{efe-swarm-sdk/} \code{src/backend/mod.rs:232} &
garde G2 \textbf{+} appariement d'épingle TLS \textbf{+} barreau d'attestation &
\best{durci} --- miroir déclaré « exact », vérifié identique ligne à ligne \\
\code{efe-swarm-gateway/} \code{src/frontier.rs:159} &
avertissement de démarrage sur l'URL du SLM &
\pire{non durci} : ne coupe pas \code{\#}, et surtout \code{starts\_with("127.")} \\
\code{efe-meta-swarm-server/} \code{controller/iterate.rs:687} &
le \textbf{SAN} attendu de la vérification mTLS &
\textcolor{amber}{\textbf{non durci}} --- aucun traitement du \code{userinfo}, ne coupe pas \code{\#} \\
\code{efe-meta-swarm-server/} \code{tls/mod.rs:145} &
\textbf{si TLS peut être coupé} sur ce lien &
\textcolor{amber}{\textbf{non durci}} --- mais sauvé par un \code{IpAddr::from\_str} strict en aval \\
\bottomrule
\end{tabular}
\end{center}

\begin{alerte}
\textbf{Le critère de fin est jouable.} Le SDK n'a pas besoin de \code{is\_local\_endpoint} seulement :
il appelle \code{host\_of} \emph{directement} pour apparier une épingle TLS
(\code{matched\_pin}) et pour estampiller l'identité d'endpoint dans un record scellé
(\code{attestation\_for}). On peut donc extraire \code{is\_local\_endpoint} dans un crate commun,
faire renvoyer \textbf{un} fichier au \code{grep} --- et laisser \code{host\_of} dupliqué, avec les
trois autres analyseurs intacts. Le module serait déclaré fini, la commande de la roadmap
sortirait 0, et \emph{le différentiel d'analyseur serait toujours là}. C'est le premier livrable de
M1 : un critère qu'on ne peut pas satisfaire sans avoir fait le travail.
\end{alerte}

\section{2 · Trois devoirs pour un seul prédicat}

Avant de décider où mettre le code, il faut voir ce qu'on lui demande. « Cet endpoint est-il local ? »
sert dans le dépôt à trois choses qui n'ont ni la même gravité, ni le même mode d'échec, ni le même
lecteur --- et c'est la raison pour laquelle la duplication a paru anodine.

\subsection{2.1 · Le devoir d'\emph{application}}

\code{assert\_local} refuse. La console de conformité l'appelle sur son adresse d'écoute
(\code{audit-console/src/server.rs:27}) : refuser de \emph{binder} ailleurs qu'en local est ce qui
rend crédible la phrase « la console ne s'expose pas ». Le SDK l'appelle sur son gate G2. C'est le
seul des trois devoirs que tout le monde reconnaît comme de la sécurité, et c'est celui que les deux
copies durcies servent correctement aujourd'hui.

\subsection{2.2 · Le devoir d'\emph{attestation} --- celui qu'on ne voit pas}

\code{efe-swarm-gateway/src/mission.rs:139} et \code{efe-swarm-sdk::attestation\_for} n'interdisent
rien. Ils \textbf{classent} : localité \fl{} \code{ComputeAttestation::R2LocalWitnessed}, distance
\fl{} \code{R4RemoteChannel}. Cette classe est ensuite \emph{scellée dans le record de provenance}.

La différence de nature est décisive et mérite d'être dite lentement. Un garde qui se trompe laisse
passer un appel ; c'est grave, c'est réparable, et cela ne laisse pas de trace mensongère. Un
\emph{attesteur} qui se trompe écrit dans un reçu non répudiable la phrase \emph{« le récepteur de
provenance était co-localisé avec l'inférence, il témoigne de la computation »} à propos d'un appel
parti à l'autre bout du monde. Le reçu est le produit. Un reçu qui ment sur la localité ne se répare
pas : il se re-prouve parfaitement, avec sa fausse affirmation dedans, et sa signature garantit
seulement que \emph{personne ne l'a modifié depuis}. C'est le pire résultat possible pour une pile
dont l'argument de vente est l'intégrité.

\begin{these}
Le devoir d'attestation transforme un faux positif de parseur en \textbf{fausse déclaration scellée}.
C'est l'argument le plus fort pour l'unicité, et c'est celui que personne n'invoque, parce que le
site d'appel ne ressemble pas à un contrôle de sécurité --- c'est un \code{match} qui remplit un
champ.
\end{these}

\subsection{2.3 · Le devoir de \emph{conseil}}

\code{frontier.rs:128} n'interdit ni ne classe : il \emph{prévient}. Si l'URL du SLM ne semble pas
locale, un avertissement part sur \code{stderr}. Le commentaire qui l'entoure est d'une honnêteté
remarquable et vaut d'être cité : \emph{« l'URL n'est PAS contrainte au local --- un llama-server
distant est possible et son egress n'est ni gardé par l'allowlist tenant ni décompté. On prévient au
boot plutôt que de prétendre une souveraineté qu'on ne vérifie pas. »} La conception est juste : on
ne ment pas, on avertit. C'est le \emph{prédicat} qui est faux --- et c'est l'objet de la section
suivante.

\section{3 · Le défaut vivant}

\subsection{3.1 · La chaîne, vérifiée bout en bout}

\code{frontier.rs:159} conclut à la localité par \code{host.starts\_with("127.")}. Or
\code{guard.rs:41} --- dans le même dépôt, dans le crate que ce même fichier importe déjà ---
documente précisément ce test comme celui contre lequel il a été durci : \emph{« Durci vs un naïf
\code{starts\_with("127.")} : \code{127.evil.com} (un NOM, qui pourrait résoudre vers l'externe) ne
doit PAS passer la frontière de souveraineté. »}

\begin{extrait}
\code{EFE\_SLM\_REFORMULATOR\_URL=http://127.evil.com:8091} \\[3pt]
\code{\ \ slm\_url\_looks\_local} \fl{} \code{host = "127.evil.com"} \fl{} \code{starts\_with("127.")} \fl{} \pire{true} \\
\code{\ \ } \fl{} aucun avertissement de démarrage \\
\code{\ \ } \fl{} le client SLM est câblé (l'URL n'est \emph{délibérément} pas contrainte) \\
\code{\ \ } \fl{} l'egress n'est ni gardé par l'allowlist tenant, ni décompté (hop budgété 0 ¢) \\
\code{\ \ } \fl{} \textbf{M0 imprime} \code{relais SLM \ \ SERVIE} \\
\code{\ \ } \fl{} \pire{chaque prompt reformulé sort du périmètre, en silence, sur un rapport vert}
\end{extrait}

Chaque maillon a été vérifié dans le code. \code{127.evil.com} est un nom DNS ordinaire, que
n'importe qui peut faire pointer où il veut ; le préfixe \code{127.} n'a aucune signification pour un
résolveur. Le test qui accompagne la fonction
(\code{slm\_url\_locality\_boot\_warn\_heuristic}) éprouve \code{10.0.0.5},
\code{slm.example.com} et \code{192.168.1.10} --- jamais un nom préfixé \code{127.}. Le défaut n'est
donc figé par aucun test : \textbf{le corriger ne casse rien}.

\subsection{3.2 · La phrase qui condamne la copie}

Le commentaire de \code{frontier.rs} justifie ainsi la troisième copie : \emph{« Le check est un
miroir minimal des \code{is\_local\_endpoint} maison --- provenance-capture/guard.rs \&
efe-swarm-sdk --- \textbf{sans en dépendre}. »} L'auteur savait donc qu'il écrivait une troisième
copie, et l'a voulue « minimale » pour éviter une dépendance.

\begin{alerte}
\textbf{La dépendance était déjà payée.} \code{frontier.rs} n'est compilé que sous la feature
\code{frontier-llm} (\code{lib.rs:58}), et c'est exactement cette feature qui active
\code{provenance-capture} dans le \code{Cargo.toml} du crate. À l'intérieur de ce fichier,
\code{provenance\_capture::is\_local\_endpoint} est \emph{déjà dans la portée} --- le fichier voisin
\code{mission.rs}, sous la même feature, l'appelle. Le « miroir minimal » n'a donc économisé
\best{aucune} dépendance. Il n'a acheté que la perte du durcissement.
\end{alerte}

Et l'ironie se referme une ligne plus loin. Le \code{Cargo.toml} de la passerelle annote ainsi la
dépendance qu'elle prend :

\begin{extrait}
\code{provenance-capture = \{ path = "../provenance-capture", optional = true \}} \\
\code{\ \ \ \ \sz{\# egress guard is\_local\_endpoint (réutilisé --- \textbf{pas de 3e copie})}}
\end{extrait}

L'annotation qui affirme qu'il n'y a pas de troisième copie et la troisième copie voyagent dans le
même crate, sous la même feature, dans le même commit.

\section{4 · La duplication était un choix --- et pourquoi il ne tient plus}

Il serait faux de présenter M1 comme la correction d'un oubli. Les deux copies durcies sont
délibérées, documentées, et défendues par un argument sérieux. \code{guard.rs:5} :
\emph{« la même règle, ré-implémentée ici pour garder la couche de capture autonome »}. Le miroir du
SDK dit la même chose, et ajoute la discipline qui devait tenir l'ensemble : \emph{« discipline
\#4 : les deux copies durcies ENSEMBLE »}.

L'argument d'autonomie n'est pas décoratif. \code{provenance-capture} ne dépend que de
\code{provenance-dag} --- une seule arête. C'est cette maigreur qui fait de la pile preuve un actif
propre : elle s'observe sur n'importe quel système d'IA sans traîner le moteur du dépôt derrière
elle. \code{efe-swarm-sdk}, lui, tire cinq crates d'essaim. Les deux fermetures sont \emph{disjointes},
et elles appartiennent à deux périmètres produit que l'analyse de rachetabilité a explicitement
demandé de ne pas souder. Faire dépendre l'un de l'autre pour économiser soixante lignes serait un
mauvais échange.

\begin{these}
Le vrai reproche à faire à la duplication n'est pas d'exister, c'est de reposer sur une
\textbf{vigilance}. « Durcir les deux ensemble » est une règle que des humains doivent se rappeler.
Elle a tenu deux fois --- les deux copies sont aujourd'hui identiques, l'audit l'a vérifié ligne à
ligne --- puis elle a échoué trois fois, dont une dans le crate qui l'avait écrite. M1 remplace une
discipline par une structure, exactement comme M0 a remplacé « pensez à documenter vos variables »
par un registre dont un test interdit de sortir.
\end{these}

\section{5 · L'architecture : un analyseur, un type, un corpus}

\subsection{5.1 · Un crate feuille, à zéro dépendance}

La contrainte de périmètre se satisfait sans compromis, parce que l'analyseur n'a besoin de rien :
\code{guard.rs} le dit lui-même, \emph{« minimal, sans dépendance (pas de crate \code{url}) »}. Un
crate \code{scella-authority} sans aucune dépendance est une \textbf{feuille} : chacun des deux
périmètres en prend une, et aucun des deux ne gagne de chemin vers l'autre. Le graphe grossit d'un
nœud ; il ne se soude pas.

Le précédent existe déjà et n'a coûté à personne : \code{audit-api/src/lib.rs:72} réexporte
\code{assert\_local}, \code{is\_local\_endpoint} et \code{SovereigntyError} depuis
\code{provenance-capture}, et c'est par cette chaîne que la console obtient son garde de \emph{bind}.
Le dépôt sait donc déjà partager ce garde --- il ne l'a simplement pas fait par-dessus la frontière
de périmètre.

\subsection{5.2 · Pourquoi un \emph{type}, et pas une chaîne}

Le crate n'exporte pas une fonction \code{host\_of} rendant \code{Option<String>}. Il exporte un type :

\begin{extrait}
\code{pub struct Authority(String);\ \ \ \sz{// constructible UNIQUEMENT par authority\_of}} \\
\code{pub fn authority\_of(endpoint: \&str) -> Option<Authority>} \\
\code{impl Authority \{ pub fn host(\&self) -> \&str ; pub fn is\_loopback(\&self) -> bool \}} \\[3pt]
\code{pub fn is\_local\_endpoint(endpoint: \&str) -> bool \ \sz{// = authority\_of(..).is\_loopback()}} \\
\code{pub fn assert\_local(endpoint: \&str) -> Result<(), SovereigntyError>}
\end{extrait}

La raison est structurelle, pas esthétique. Toute la classe de bugs se résume à : \emph{quelqu'un a
re-dérivé un hôte depuis une chaîne}. Un \code{String} nu invite à la re-dérivation --- on le
concatène, on le recoupe, on le compare à un autre hôte obtenu autrement. Un \code{Authority} que
seul \code{authority\_of} peut fabriquer rend l'accident coûteux : pour obtenir un hôte, il faut
passer par l'analyseur. C'est ce dont \code{matched\_pin} a besoin --- comparer l'hôte d'un endpoint
à celui d'une épingle exige que les \emph{deux} viennent du même analyseur, ce que le type impose et
qu'une fonction rendant \code{String} laisserait au hasard.

\subsection{5.3 · Le corpus de conformité, exporté}

Le crate exporte enfin, en \code{pub const}, la table des cas qui l'ont durci : les six bypass SSRF,
les formes normalisées, les autorités \emph{fail-close}. Tout consommateur qui, pour une raison
légitime, garderait un analyseur propre peut alors lui faire passer \emph{le même} corpus. C'est le
filet pour les cas où l'extraction est impossible --- et c'est ce qui aurait attrapé les copies
\#4 et \#5, que personne n'a jamais confrontées aux cas du garde.

\section{6 · La question de fond : à la main, ou sur \code{url} ?}

C'est la décision structurante de M1, et elle mérite d'être posée franchement plutôt que héritée du
statu quo. L'analyseur cherche à répondre à une question qui n'est pas « que dit la RFC 3986 ? » mais
\textbf{« quel hôte \code{reqwest} va-t-il réellement contacter ? »}. Ce sont deux questions
différentes, et tout l'écart entre elles est de la surface d'attaque.

\begin{center}\footnotesize
\renewcommand{\arraystretch}{1.25}
\begin{tabular}{@{}>{\raggedright\arraybackslash}p{3.1cm} >{\raggedright\arraybackslash}p{6.2cm} >{\raggedright\arraybackslash}p{6.0cm}@{}}
\toprule
\textbf{Option} & \textbf{Ce qu'elle donne} & \textbf{Ce qu'elle coûte} \\
\midrule
\textbf{A} --- à la main, zéro dépendance \sz{(statu quo extrait)} &
crate feuille, auditable en un écran, aucun périmètre alourdi &
\pire{la liste \emph{fail-close}} (\code{\textbackslash}, espace, contrôle) est une \textbf{liste noire} --- elle n'est jamais complète, et a déjà exigé deux tours de durcissement \\
\textbf{B} --- bâti sur \code{url} &
\best{accord avec \code{reqwest} par construction} : la classe entière disparaît, pas seulement ses instances connues &
\code{url} tire \code{idna}, \code{percent-encoding}, \code{unicode-*} \fl{} un arbre réel dans une pile preuve qui n'a \textbf{qu'une} arête ; et l'IDNA ouvre sa propre surface (homographes punycode) \\
\textbf{C} --- à la main, \best{prouvé} contre \code{url} en \code{dev-dependencies} &
le crate livré reste une feuille, \emph{et} la liste noire devient une \textbf{propriété vérifiée} &
le test ne vaut que son corpus --- mais un corpus engendré couvre incomparablement plus que six cas écrits à la main \\
\bottomrule
\end{tabular}
\end{center}

\begin{acquis}
\textbf{Décision : C.} Le crate reste sans dépendance en production, et un test de propriété --- avec
\code{url} en \code{dev-dependencies} seulement --- vérifie sur un corpus engendré l'invariant
suivant, qui est la formulation exacte de ce que le garde promet :

\smallskip
\code{pour tout endpoint \emph{x} : \ authority\_of(x) == None \ \ \textbf{OU} \ \ authority\_of(x).host() == Url::parse(x).host\_str()}
\smallskip

Autrement dit : \emph{ou bien nous refusons, ou bien nous désignons exactement l'hôte que le client
appellera}. Jamais un troisième cas. C'est le même geste que M0 : remplacer « nous avons pensé à ces
attaques » par « nous l'avons prouvé, ou nous refusons ».
\end{acquis}

\begin{alerte}
\textbf{Vérifié à l'implémentation (M1.a) : la propriété a trouvé \pire{deux} différentiels réels
dès sa première exécution hostile} --- deux que ni la rédaction de cette architecture, ni deux tours
de durcissement, ni aucune revue n'avaient vus.

\smallskip
\code{http://2130706433} : \code{url} --- donc \code{reqwest} --- normalise l'entier décimal en
\code{127.0.0.1} ; l'analyseur rendait \code{2130706433}. Même famille pour \code{0177.0.0.1} (le
zéro initial vaut octal pour un résolveur historique) et \code{127.1} (forme courte).

\smallskip
\code{http://[::0001]} : \code{url} canonise en \code{::1} ; l'analyseur rendait \code{::0001}. Et
en instrumentant, un troisième cas est apparu, présent dans les \emph{deux} implémentations
historiques : \code{http://fe80::1:8080} --- IPv6 nue --- produisait l'hôte \code{fe80}, une chaîne
qui ne désigne rien.

\smallskip
Les trois étaient \best{sûrs pour la localité} : « non local » dans tous les cas, donc jamais un
laissez-passer. Mais \code{host\_of} sert aussi à \textbf{apparier une épingle TLS} et à
\textbf{nommer un endpoint dans un journal} : désigner un hôte différent de celui appelé y a des
conséquences propres. La règle retenue n'est pas d'imiter ces normalisations --- les imiter
reviendrait à parier sur celle qu'applique le client --- mais de \textbf{refuser toute écriture
ambiguë}. Le crate compte donc \textbf{six} cas de refus, non quatre.

\smallskip
C'est l'argument de l'option C, réalisé : une liste noire contient ce à quoi quelqu'un a pensé ;
une propriété contient ce qui est vrai.
\end{alerte}

\section{7 · La migration, site par site}

\begin{center}\footnotesize
\renewcommand{\arraystretch}{1.25}
\begin{tabular}{@{}>{\raggedright\arraybackslash}p{3.6cm} >{\raggedright\arraybackslash}p{6.3cm} >{\raggedright\arraybackslash}p{5.4cm}@{}}
\toprule
\textbf{Site} & \textbf{Geste} & \textbf{Changement de comportement} \\
\midrule
\code{provenance-capture/} \code{guard.rs} &
devient une \textbf{réexportation} du crate feuille &
\best{aucun} --- et les 5 dépendants (\code{audit-api}, \code{audit-console}, \code{reg-evidence}, \code{forensic-export}, passerelle) ne bougent pas \\
\code{efe-swarm-sdk/} \code{backend/mod.rs} &
supprime \code{host\_of} et \code{is\_local\_endpoint} ; \code{matched\_pin} compare des \code{Authority} &
\best{aucun} --- les deux copies étaient déjà identiques \\
\code{efe-swarm-gateway/} \code{frontier.rs} &
supprime \code{slm\_url\_looks\_local}, appelle le prédicat \textbf{déjà dans la portée} &
\textcolor{amber}{\textbf{oui, et c'est le but}} : \code{127.evil.com} avertit désormais ; \code{LOCALHOST} cesse d'avertir à tort \\
\code{meta-swarm/} \code{iterate.rs} \sz{(SAN)} &
\code{authority\_of} pour l'\emph{analyse}, décision de SAN inchangée &
à mesurer : les entrées malformées tombaient déjà sur le SAN configuré \\
\code{meta-swarm/} \code{tls/mod.rs} &
\code{authority\_of} pour l'analyse, \code{IpAddr} conservé pour l'\emph{acceptation} &
\pire{oui} --- voir ci-dessous \\
\bottomrule
\end{tabular}
\end{center}

\begin{alerte}
\textbf{Le seul risque de régression, déclaré.} \code{tls/mod.rs} accepte aujourd'hui une IPv6 nue
sans crochets (\code{authority.contains("::")} \fl{} l'autorité entière est traitée comme l'hôte).
L'analyseur durci, lui, \emph{refuse} \code{http://::1:8080} : sans crochets, la coupe sur \code{:}
rend un hôte vide, donc \code{None}, donc « non local ». Or dans ce fichier, « non local » veut dire
\best{TLS reste allumé} --- la direction sûre. Un déploiement qui écrit ses endpoints en IPv6 nue
verrait donc du TLS apparaître là où il n'y en avait pas. C'est un durcissement, pas une panne, mais
c'est un changement observable en production : il exige son propre test, sa propre note de version,
et c'est la raison pour laquelle ces deux sites sont sortis du premier incrément.

\smallskip
\textbf{Corrigé après M1.a : la portée est plus large qu'écrit ici.} L'analyseur partagé refuse
\emph{trois} familles et non une --- IPv6 nue, IPv6 non canonique (\code{[::0001]}), et écritures
numériques ambiguës d'IPv4. Un déploiement qui emploie l'une d'elles verra donc TLS s'allumer là où
il était coupé. Toujours un durcissement, toujours la direction sûre, mais la note de version de
M1.b doit nommer les trois.
\end{alerte}

Il faut être juste envers ces deux fichiers : ils sont \emph{sûrs aujourd'hui}. Les charges de
différentiel (\code{\#@}, \code{userinfo@}) leur font produire un hôte qui n'est ni
\code{localhost} ni une IP analysable, donc \code{false}, donc TLS maintenu et SAN inutilisable. Mais
cette sûreté vient de la \emph{sévérité de ce qui suit} --- \code{IpAddr::from\_str}, l'égalité
stricte de chaîne, l'absence de certificat correspondant --- et non de l'analyseur. Ils sont sûrs par
composition, pas par conception : ce sont eux qui casseront le jour où quelqu'un les
« améliorera ».

\section{8 · Le critère d'acceptation --- exécutable et non jouable}

Le critère de la roadmap est conservé mais \textbf{subordonné} : il reste vrai, il ne suffit plus.

\begin{extrait}
\code{\sz{1.}\ grep -rl '"://"' crates/*/src \ \fl{} \textbf{1} fichier \ \sz{(5 aujourd'hui)}} \\
\code{\sz{2.}\ cargo test -p scella-authority --test un\_seul\_analyseur \ \fl{} vert} \\
\code{\sz{3.}\ cargo test -p scella-authority --test accord\_avec\_url \ \fl{} vert} \\
\code{\sz{4.}\ EFE\_SLM\_REFORMULATOR\_URL=http://127.evil.com cargo run ... 2>\&1 | grep 'NON garantie'}
\end{extrait}

Le premier compte les \emph{analyseurs}, pas les gardes : c'est lui qui rend le critère non jouable,
et c'est lui qui aurait signalé les copies \#3, \#4 et \#5 le jour où elles ont été écrites.

\begin{alerte}
\textbf{Ce critère a lui-même dû être mesuré, et sa première rédaction était fausse.} Elle cherchait
\code{split\_once("://")} --- la forme employée par quatre des cinq analyseurs. Le cinquième,
\code{frontier.rs}, écrit \code{url.split("://").nth(1)}. Le grep aurait donc renvoyé \textbf{quatre}
fichiers, serait tombé à un après extraction des quatre, aurait déclaré M1 fini --- et aurait laissé
en place \emph{précisément la copie qui porte le défaut vivant}. C'est le même piège que le critère
de la roadmap, une couche plus bas. Le motif retenu est le littéral \code{"://"} seul : mesuré, il
désigne exactement les cinq analyseurs et rien d'autre. \best{Un critère qui n'a pas été exécuté
n'est pas un critère.}
\end{alerte} Le
deuxième est l'analogue direct du \code{T1} de M0 --- un test qui parcourt l'arbre des sources et
échoue sur toute occurrence hors du crate, avec une liste d'exemptions qui doit rester
\textbf{vide}. Le troisième est la propriété d'accord de la section 6. Le quatrième est le défaut
vivant : il \emph{échoue aujourd'hui}, et c'est la meilleure garantie qu'il mesure quelque chose.

\begin{these}
Un critère de fin qu'on peut atteindre sans faire le travail n'est pas un critère, c'est une
formalité. Celui-ci a été écrit \emph{après} l'inventaire, ce qui est la seule façon d'obtenir un
critère qui garde un fait plutôt qu'une hypothèse.
\end{these}

\section{9 · Le périmètre --- ce que M1 ne fait pas}

\textbf{M1 ne contraint pas l'URL du SLM.} La liberté d'un \code{llama-server} distant est un choix
assumé et documenté ; M1 corrige l'\emph{avertissement}, il ne ferme pas la porte. Contraindre serait
un changement de produit, pas d'architecture, et il appartiendrait à M0 (un interrupteur, sa classe,
son coût d'absence) ou à la voie C3.

\textbf{M1 n'ajoute pas de résolution DNS.} Le garde raisonne sur le \emph{littéral}. Un nom qui
résout vers \code{127.0.0.1} n'est pas « local » pour lui, et c'est correct : résoudre au moment du
contrôle ouvrirait un TOCTOU (la réponse DNS peut changer entre le contrôle et l'appel), qui est un
piège plus profond que celui qu'on referme.

\textbf{M1 ne touche pas à l'allowlist d'egress par tenant} ni au comptage de budget. Ce sont des
mécanismes voisins, avec leur propre analyse ; les mêler ferait de M1 un module d'une semaine.

\textbf{M1 n'unifie pas les \emph{décisions}, seulement l'\emph{analyse}.} \code{tls/mod.rs} peut
continuer d'accepter \code{127.1.2.3} par \code{IpAddr} et le SDK d'exiger quatre octets : ce sont des
politiques distinctes, légitimement distinctes. Ce qui doit être unique, c'est la réponse à
« \emph{quel hôte} ? ». Confondre les deux serait la sur-ingénierie qui fait échouer ce genre de
module.

\section{10 · L'ordre, et la taille réelle}

La roadmap donne deux jours. L'inventaire ayant fait passer le périmètre de deux sites à cinq, dont
deux avec un changement de comportement observable, la taille honnête est de trois à quatre jours ---
sauf à découper, ce qui est ici la bonne réponse.

\begin{center}\small
\renewcommand{\arraystretch}{1.25}
\begin{tabular}{@{}l >{\raggedright\arraybackslash}p{9.3cm} >{\raggedright\arraybackslash}p{3.5cm} c@{}}
\toprule
& \textbf{Contenu} & \textbf{Fini quand} & \sz{Taille} \\
\midrule
\textbf{M1.a} \best{fait} & le crate feuille (843 l.), le type \code{Authority}, le corpus exporté ; les sites de souveraineté migrés --- \code{frontier.rs} par P1, \code{guard.rs} devenu ré-export, le SDK adopté ; le test d'unicité ; la propriété d'accord & critères 2 à 4 ; le 1\textsuperscript{er} tombe à \textbf{3} & \sz{2 j} \\
\textbf{M1.b} & les \textbf{deux} sites méta-essaim ; convergence de la règle d'acceptation ; note de version sur l'IPv6 nue & critère 1 étendu au méta-essaim & \sz{1--2 j} \\
\bottomrule
\end{tabular}
\end{center}

M1.a livre le critère de la roadmap, ferme le défaut vivant, et n'a \emph{aucun} changement de
comportement en dehors de l'avertissement qu'on voulait corriger --- ce que prouve directement le
fait que les \textbf{tests historiques} de \code{guard.rs}, laissés en place et pointés sur la
nouvelle implémentation, passent sans être touchés. La couture annoncée est tenue : le critère
\code{"://"} tombe de cinq à \textbf{trois}, pas à un, et les deux analyseurs restants sont
\textbf{exemptés nommément}, avec leur raison et le module qui les supprimera --- une assertion
interdit à cette liste de grandir, et le critère de fin de M1.b est qu'elle soit vide. M1.b porte tout le risque et
aucune urgence : ses deux sites sont sûrs aujourd'hui. Les séparer n'est pas un report, c'est mettre
le risque là où il se voit.

\begin{acquis}
\textbf{Ce que M1 achète, en une phrase.} Aujourd'hui, la souveraineté du produit --- « vos données
ne sortent pas » --- repose sur cinq réponses indépendantes à une même question, dont trois sont
fausses ou fragiles, et sur une discipline humaine qui a déjà échoué trois fois. Après M1, elle
repose sur \textbf{un} analyseur, dont un test interdit la copie et dont une propriété prouve
l'accord avec le client HTTP réel. Le gain n'est pas soixante lignes de moins : c'est que la question
« et si les deux copies divergeaient ? » \best{cesse d'être posable}.
\end{acquis}

\end{document}
